\documentclass[11pt,a4paper]{article}

% Idioma y fuentes (compilar con XeLaTeX)
\usepackage{fontspec}
\usepackage[spanish]{babel}

\usepackage[margin=2.5cm]{geometry}
\usepackage{amsmath, amssymb}
\usepackage{graphicx}
\usepackage{enumitem}
\usepackage[hidelinks]{hyperref}

% Datos de la tarea
\newcommand{\curso}{Nombre del curso}
\newcommand{\tarea}{Tarea 1}
\newcommand{\estudiante}{Nombre Apellido}
\newcommand{\profesor}{Nombre del profesor}

% Pregunta numerada y entorno de respuesta
\newcounter{pregunta}
\newcommand{\pregunta}[1]{%
  \stepcounter{pregunta}%
  \section*{Pregunta \thepregunta\ifx\relax#1\relax\else: #1\fi}}
\newenvironment{respuesta}{\par\noindent\textbf{Respuesta.}\ }{\hfill$\blacksquare$\par}

\begin{document}

\noindent
\begin{minipage}[t]{0.6\textwidth}
  \textbf{\curso}\\
  \tarea
\end{minipage}%
\begin{minipage}[t]{0.4\textwidth}
  \raggedleft
  \estudiante\\
  Profesor: \profesor\\
  \today
\end{minipage}
\par\medskip\hrule\medskip

\pregunta{Derivadas}
Calcule $\dfrac{d}{dx}\left(x^2 \sin x\right)$.

\begin{respuesta}
  Por la regla del producto,
  \[
    \frac{d}{dx}\left(x^2 \sin x\right) = 2x \sin x + x^2 \cos x.
  \]
\end{respuesta}

\pregunta{}
Enuncie dos propiedades de la integral definida.

\begin{respuesta}
  \begin{enumerate}[label=(\alph*)]
    \item Linealidad: $\int_a^b (\alpha f + \beta g) = \alpha \int_a^b f + \beta \int_a^b g$.
    \item Aditividad: $\int_a^b f = \int_a^c f + \int_c^b f$.
  \end{enumerate}
\end{respuesta}

\end{document}
